\section{内网穿透}
\begin{code}[title=Rathole Service]
	\begin{verbatim}
sudo vim /etc/systemd/system/ratholes.service

[Unit]
Description=Rathole Server Service
After=network.target

[Service]
Type=simple
Restart=on-failure
RestartSec=5s
LimitNOFILE=1048576

# with root
ExecStart=/usr/bin/rathole -s /etc/rathole/rathole.toml
# without root
# ExecStart=%h/.local/bin/rathole -s %h/.local/etc/rathole/rathole.toml

[Install]
WantedBy=multi-user.target
    \end{verbatim}
\end{code}

\begin{code}[title=Rathole Toml]
	\begin{verbatim}
sudo mkdir -p /etc/rathole/rathole.toml

# server.toml
[server]
bind_addr = "0.0.0.0:2333" # `2333` 配置了服务端监听客户端连接的端口

[server.services.my_nas_ssh]
token = "use_a_secret_that_only_you_know" # 用于验证的 token
bind_addr = "0.0.0.0:5202" # `5202` 配置了将 `my_nas_ssh` 暴露给互联网的端口


sudo systemctl daemon-reload
sudo systemctl enable ratholes
sudo systemctl start  ratholes


# client.toml
[client]
remote_addr = "myserver.com:2333" # 服务器的地址。端口必须与 `server.bind_addr` 中的端口相同。
[client.services.my_nas_ssh]
token = "use_a_secret_that_only_you_know" # 必须与服务器相同以通过验证
local_addr = "127.0.0.1:22" # 需要被转发的服务的地址
    \end{verbatim}
\end{code}
